\begin{lemma}[Type filters preserve substitutability]
\label{lem:type-filter}
If $L$ is $(f,h)$-tuple-substitutable and $F\subseteq M$, then
$L\cap h^{-1}(F)$ is $(f,h)$-tuple-substitutable.
\end{lemma}

\begin{proof}
Let $\tuple x,\tuple y$ have equal $h$-vectors and share a context $E$
in $L':=L\cap h^{-1}(F)$.  They share $E$ in $L$, so they have the same
$L$-distribution.  If $E'[\tuple x]\in L'$, then $E'[\tuple y]\in L$
and $h(E'[\tuple y])=h(E'[\tuple x])\in F$, so $E'[\tuple y]\in L'$.
\end{proof}

\begin{corollary}[Pseudovariety dichotomy]
\label{cor:pseudovariety-dichotomy}
Let $\mathbf V$ be a pseudovariety of finite semigroups.
\begin{enumerate}[label=\textup{(\roman*)},leftmargin=*]
\item If $\mathbf V\subseteq\mathbf{LG}$, then for every alphabet,
every typing $h$ with $h(\Sigma^+)\in\mathbf V$, and all fixed
$f,r_{\max}$, characteristic data of total length
$O(\|G\|^3\bar\theta_G)$ exist as in
Theorem~\ref{thm:local-group-polydata}.
\item If $\mathbf V\not\subseteq\mathbf{LG}$, then
$U_1\in\mathbf V$.  At fan-out one there is a typing $h_{\mathrm U}$ with
$h_{\mathrm U}(\Sigma^+)=U_1$ and targets with $\|G\|,\theta_G=O(n)$,
every characteristic sample of which for
$\mathcal B^{\le2}_{1,h_{\mathrm U}}$ contains a word of length
$3\cdot2^n-2$.  The same holds for every set-driven operator admitting
characteristic samples throughout $\Cmcf{1}{h_{\mathrm U}}$.
\end{enumerate}
\end{corollary}

\begin{proof}
(i) is Theorem~\ref{thm:local-group-polydata}.  For (ii), choose
$S\in\mathbf V\setminus\mathbf{LG}$.  By
Theorem~\ref{thm:lg-characterization}, $S$ has a subsemigroup isomorphic
to $U_1$, and pseudovarieties are closed under subsemigroups.

Take the doubling family $L_n$, the typing $h_{\mathrm U}$ and the long
type-$1$ word $w_n^*$ of~\cite[Proposition~6.3]{kuriyamaCFG}.  There
$L_n$ is shown to be Clark--Eyraud substitutable, hence
$(1,h_{\mathrm U})$-tuple-substitutable.  By
Lemma~\ref{lem:type-filter},
$L_n^-:=L_n\cap h_{\mathrm U}^{-1}(\zeta)=L_n\setminus\{w_n^*\}$ is as
well.  If $K\subseteq L_n$ omits $w_n^*$, then $K\subseteq L_n^-$, and
Proposition~\ref{prop:reconstruction-soundness} gives
$L(\mathcal B^{\le2}_{1,h_{\mathrm U}}(K))\subseteq L_n^-$.  Thus every
characteristic sample contains $w_n^*$.  The set-driven statement is
the distinguishing-pair argument of~\cite[Proposition~6.3(iv)]{kuriyamaCFG} for the pair
$L_n,L_n^-$.
\end{proof}

\begin{remark}[Scope of the dichotomy]
\label{rem:dichotomy-scope}
\begin{enumerate}[label=\textup{(\roman*)},leftmargin=*]
\item Part~(i) holds at every fixed fan-out.  Part~(ii) is established
only at fan-out one.  The doubling family does not lift directly: $L_2$
is not $(2,h_{\mathrm U})$-tuple-substitutable.  The type-$1$ pairs
$(\mathtt[,\mathtt a\mathtt])$ and $(\mathtt[\mathtt a,\mathtt])$ share
the context $\mathtt{[[aa]}\,\blank_1\,\mathtt a\,\blank_2\,\mathtt]$
but have different distributions.  This was checked exhaustively for
$n\le3$.  Under the finer typing recording each component's set of
letters, which still has semilattice image, the corresponding failure
appears at $n=3$.  Whether part~(ii) holds
at fan-out $f\ge2$ is the pseudovariety form of
Open Problem~\ref{op:fanout-two}.
\item The dichotomy concerns pseudovarieties.  For an individual typing
$h$ with $h(\Sigma^+)\notin\mathbf{LG}$, the class $\Cmcf{1}{h}$ need
not contain a coded copy of the doubling family.  Whether every such
typing forces superpolynomial characteristic data remains open.
\end{enumerate}
\end{remark}
